\section{Introducción}

El gasto en salud representa uno de los principales desafíos financieros
para individuos y sistemas de cobertura médica. En Estados Unidos, el gasto
per cápita en salud superó los \$12,500 dólares anuales en 2020,
representando el 19.7\% del PIB \cite{cms2022national}. La capacidad de
predecir con precisión los costos médicos individuales permite a las
aseguradoras fijar primas justas, gestionar riesgos y diseñar programas
de prevención dirigidos.

Los modelos de regresión supervisada han demostrado ser herramientas
eficaces para la predicción de costos médicos, capturando tanto relaciones
lineales como interacciones complejas entre variables demográficas y
factores de riesgo como el tabaquismo o el Índice de Masa Corporal
(IMC) \cite{duan2021prediction}.

Este trabajo implementa y compara múltiples algoritmos de regresión sobre
el Medical Cost Personal Dataset, explorando el impacto de la ingeniería
de features en el rendimiento predictivo. Se presta especial atención a la
interacción BMI-tabaquismo, identificada en la literatura como uno de los
predictores más potentes del gasto médico elevado.

\subsection{Objetivos}
\begin{itemize}
  \item Comparar seis modelos de regresión aplicando CRISP-DM.
  \item Evaluar el impacto de la ingeniería de features (interacciones,
        transformaciones no lineales) en la predicción.
  \item Identificar los factores más influyentes en el costo del seguro.
  \item Desplegar el modelo ganador como servicio web con Flask y Plotly.
\end{itemize}
